\documentclass[10pt,a4paper]{amsart}
\usepackage[margin=19mm]{geometry}
\usepackage{amsmath,amssymb,mathtools}
\usepackage[scr=rsfs,cal=euler]{mathalfa}
\usepackage{xfrac}
\usepackage[unicode,hidelinks,bookmarks=false]{hyperref}
\newcommand{\ie}{i.e.\ }
\newcommand{\cf}{cf.\ }
\newcommand{\resp}{respectively\ }
\newcommand{\Subsection}{\subsection}
\providecommand{\nobreakdash}{\nobreak\hbox{-}}
\setlength{\emergencystretch}{2em}
\multlinegap0pt
\usepackage{bm}
\usepackage{amssymb}
\usepackage{xcolor}
\usepackage{algorithm}\usepackage{algpseudocode}
\newtheorem{definition}{Definition}
\newtheorem{theorem}{Theorem}
\newcommand{\fnc}[1]{\ensuremath{\mathcal{#1}}}
\newcommand{\mat}[1]{\ensuremath{\mathsf{#1}}}
\title{A collocation scheme that is equivalent to discontinuous Galerkin discretizations}
\author{Jason Hicken}
\date{Source: arXiv:2605.27327v1 (CC BY 4.0)}
\begin{document}
\maketitle

\noindent\emph{Source and licence.} A collocation scheme that is equivalent to discontinuous Galerkin discretizations. Version arXiv:2605.27327v1 (CC BY 4.0), distributed by arXiv under the Creative Commons Attribution 4.0 International licence, http://creativecommons.org/licenses/by/4.0/, as recorded in the arXiv licence metadata for this exact version and shown on the abstract page https://arxiv.org/abs/2605.27327v1. Full licence notice: \texttt{licenses/arxiv-2605.27327-CC-BY-4.0.txt}.\par\smallskip

\noindent\textit{Editorial scope.} Counted unit: the article's theorem thm:mc\_are\_sbp (source0.tex lines 231--233). The article's complete original proof of it is delivered with it (source0.tex lines 235--262, one \texttt{proof} environment of 28 lines). No mathematics is written, repaired or re-derived: the statement and the proof are the article's own lines, in the article's own notation, and no retained line is substituted or re-flowed.  Retained uncounted as the source-local context the proof needs: lines 178--181; lines 187--191; lines 196--211.  Nothing was omitted from any retained span, so the package carries no omission record.  No retained line contains a citation, so the wrapper carries no bibliography and no bibliography entry.  No retained line contains a figure environment, an image-inclusion command or drawing code, so no image is delivered and the package declares no asset dependency.  Editorial formatting only, in addition: an AMS \texttt{amsart} preamble that restates the article's own declarations and macros used by the retained lines, namely the environments \texttt{definition}, \texttt{theorem} on separate counters, together with the standard packages those lines need.  The retained environments print in this document as Definition 1, Theorem 1 in order of appearance, whereas the article prints the same items with the numbering of its own full document.  The complete native source of the article as distributed in the arXiv e-print for this exact version is bundled unchanged as \texttt{sources/201487/source0.tex}.
\begin{equation}\label{eq:ortho}
    \int_{\Omega} \fnc{V}_i \fnc{V}_j \,d\Omega = \delta_{ij},
    \qquad\forall\, i,j = 1, 2, \ldots, N_P,
\end{equation}

\begin{equation}\label{eq:vandermonde}
    [\mat{V}]_{ij} = \fnc{V}_{j}(\bm{x}_i), \qquad 
    \text{and}\qquad 
    [\mat{V}_{x_d}]_{ij} = \frac{\partial \fnc{V}_{j}}{\partial x_d}(\bm{x}_i),
\end{equation}

\begin{definition}[Diagonal-norm, first-derivative summation-by-parts operator]\label{def:sbp}
The matrix $D_{x_d} \in \mathbb{R}^{N\times N}$ is a degree $P$, diagonal-norm, summation-by-parts operator approximating the derivative operator $\partial/\partial x_d$, $d=1,2,\ldots,D$, at the nodes $X_{\Omega} = \{x_{i}\}_{i=1}^{N}$ if it satisfies the following conditions.
\begin{description}
    \item[Derivative Accuracy:] The matrix differentiates polynomials in $\mathbb{P}^P(\Omega)$ exactly.  Thus,
    \begin{equation}\label{eq:sbp_accuracy}
        \mat{D}_{x_d} \mat{V} = \mat{V}_{x_d}
    \end{equation}
    where $\mat{V}$ and $\mat{V}_{x_d}$ are the matrices defined in~\eqref{eq:vandermonde}.
    \item[Factorization:] The matrix can be factored as $\mat{D}_{x_d} = \mat{W}^{-1} \mat{Q}_{x_d}$, where $\mat{W}$ is a diagonal, positive-definite matrix.
    \item[Boundary Accuracy:] The symmetric matrix $\mat{E}_{x_d} = \mat{Q}_{x_d} + \mat{Q}_{x_d}^{T}$ satisfies
    \begin{equation}\label{eq:sbp_boundary}
    \big[ \mat{V}^T \mat{E}_{x_d} \mat{V} \big]_{ij} = \int_{\Gamma} \fnc{V}_i \fnc{V}_j \, n_{x_d}\, d\Gamma,
    \end{equation}
    where $n_{x_d}$ is the $x_{d}$ component of the outward unit normal vector on $\Gamma$.
\end{description}
\end{definition}

\begin{theorem}\label{thm:mc_are_sbp}
All degree $P$ first-derivative MC operators are also degree $P$, diagonal-norm SBP operators.
\end{theorem}

\begin{proof}
The derivative-accuracy condition follows by multiplying the MC operator and the Vandermonde matrix $\mat{V}$:
\begin{equation*}
    \mat{D}_{x_d} \mat{V} = \mat{V}_{x_d} \mat{V}^T \mat{W} \mat{V}
     = \mat{V}_{x_d},
\end{equation*}
where I used the identity $\mat{V}^T \mat{W} \mat{V} = \mat{I}_{N_P}$, which follows from the orthonormality of the basis, c.f.~\eqref{eq:ortho}, and the $2P$ exactness of the quadrature used in the MC operator definition.

The factorization condition in Definition~\ref{def:sbp} is trivially satisfied by defining 
\begin{equation}\label{eq:Q_identity}
    \mat{Q}_{x_d} \equiv \mat{W} \mat{D}_{x_d} = \mat{W} \mat{V}_{x_d} \mat{V}^T \mat{W}.
\end{equation}
This definition of $\mat{Q}_{x_d}$ implies that the symmetric matrix 
\begin{equation}\label{eq:E_identity}
    \mat{E}_{x_d} = \mat{Q}_{x_d} + \mat{Q}_{x_d}^T
    = \mat{W} \mat{V}_{x_d} \mat{V}^T \mat{W} + \mat{W} \mat{V} \mat{V}_{x_d}^T \mat{W},
\end{equation}
which can be used to verify the boundary-accuracy condition, Equation~\eqref{eq:sbp_boundary}, as follows:
\begin{alignat*}{2}
\big[ \mat{V}^T \mat{E}_{x_d} \mat{V} \big]_{ij}
&= \Big[ \mat{V}^T \Big( \mat{W} \mat{V}_{x_d} \mat{V}^T \mat{W} + \mat{W} \mat{V} \mat{V}_{x_d}^T \mat{W} \Big) \mat{V} \Big]_{ij},\qquad& &\text{(using \eqref{eq:E_identity})} \\
&= \big[ \mat{V}^T \mat{W} \mat{V}_{x_d} + \mat{V}_{x_d}^T \mat{W} \mat{V} \big]_{ij}, & &\text{(using $\mat{V}^T \mat{W} \mat{V} = \mat{I}_{N_P}$)}\\
&= \sum_{k=1}^{N} \fnc{V}_i(\bm{x}_k) w_k \frac{\partial \fnc{V}_j}{\partial x_d}(\bm{x}_k) + \sum_{k=1}^{N} \frac{\partial \fnc{V}_i}{\partial x_d} (\bm{x}_k) w_k \fnc{V}_j (\bm{x}_k) && \\
&= \int_{\Omega} \Big( \fnc{V}_i \frac{\partial \fnc{V}_j}{\partial x_d} + \frac{\partial \fnc{V}_i}{\partial x_d} \fnc{V}_{j} \Big) \, d\Omega &&\text{(using quad. accuracy)}\\
&= \int_{\Omega} \fnc{V}_i \fnc{V}_j \, n_{x_d}\, d\Gamma, &&\text{(using div. theorem)}
\end{alignat*}
as required.  Thus, all three conditions are satisfied, which concludes the proof.
\end{proof}

\end{document}
